%% DeFuzz — IEEE S&P submission draft
%% Compile: latexmk -pdf main.tex   (or: pdflatex; bibtex; pdflatex; pdflatex)
\documentclass[conference]{IEEEtran}
\IEEEoverridecommandlockouts

\usepackage{cite}
\usepackage{amsmath,amssymb,amsfonts}
\usepackage{algorithmic}
\usepackage{graphicx}
\usepackage{textcomp}
\usepackage{xcolor}
\usepackage{booktabs}
\usepackage{array}
\usepackage{listings}
\usepackage{url}
\usepackage[hidelinks]{hyperref}

\graphicspath{{figures/}}

% ---- rendered figure: academic-style D2 -> vector PDF (see figures/src/*.d2).
% usage: \incfig{stem}{caption} -> includes figures/stem.pdf, label fig:stem
\newcommand{\incfig}[2]{%
  \begin{figure}[t]\centering
  \includegraphics[width=\columnwidth]{#1}%
  \caption{#2}\label{fig:#1}
  \end{figure}%
}
% single-column figure with an explicit height cap (for tall diagrams)
\newcommand{\incfigh}[3]{%
  \begin{figure}[t]\centering
  \includegraphics[width=#3\columnwidth]{#1}%
  \caption{#2}\label{fig:#1}
  \end{figure}%
}
\newcommand{\wincfig}[2]{%
  \begin{figure*}[t]\centering
  \includegraphics[width=0.86\textwidth]{#1}%
  \caption{#2}\label{fig:#1}
  \end{figure*}%
}

% ---- placeholder figure: a labeled box inside a real float so \ref resolves.
% usage: \figplaceholder{stem}{caption}  -> label becomes fig:stem
\newcommand{\figplaceholder}[2]{%
  \begin{figure}[t]\centering
  \setlength{\fboxsep}{10pt}%
  \fbox{\parbox[c][2.2cm][c]{0.86\columnwidth}{\centering
    \textbf{[FIGURE PLACEHOLDER]}\\[3pt]\small\ttfamily #1}}%
  \caption{#2}\label{fig:#1}
  \end{figure}%
}
\newcommand{\wfigplaceholder}[2]{%
  \begin{figure*}[t]\centering
  \setlength{\fboxsep}{10pt}%
  \fbox{\parbox[c][2.6cm][c]{0.9\textwidth}{\centering
    \textbf{[FIGURE PLACEHOLDER]}\\[3pt]\small\ttfamily #1}}%
  \caption{#2}\label{fig:#1}
  \end{figure*}%
}
\newcommand{\tbd}[1]{{\color{red}\textbf{[TBD:~#1]}}}

\lstset{basicstyle=\ttfamily\footnotesize,breaklines=true,columns=fullflexible,frame=single,xleftmargin=2pt}

\begin{document}

\title{DeFuzz: Agentic Discovery of Silently-Failing\\Compiler Defenses via Cross-Mechanism Security Invariants}

\author{\IEEEauthorblockN{Anonymous Submission}
\IEEEauthorblockA{Paper \#XXX --- Under Review}}

\maketitle

\begin{abstract}
Compiler-inserted defenses---stack canaries, \texttt{\_FORTIFY\_SOURCE},
CET/IBT, CFI, BTI, PAC, and shadow stacks---are the last line that must hold
once a memory-safety bug is triggered. Their effectiveness depends not only on
the correctness of the defense code but also on assumptions about the
underlying instruction set architecture (ISA). When such an assumption breaks
on a particular backend, the defense can fail \emph{silently}: the program
compiles without warning and runs with correct functionality, yet its security
contract is no longer enforced. CVE-2023-4039, in which the stack protector of GCC
is defeated on AArch64 in the presence of variable-length arrays, is one such
case, and its cross-architecture siblings have persisted upstream for years.

The detection of silent failures is challenging due to two coupled reasons. The search space is
a two-dimensional matrix of defense mechanism $\times$ ISA, each cell backed by
an independent middle-end pass or backend template. Furthermore, failure
adjudication is difficult: the binary neither crashes nor disagrees with other
compilers; thus, crash-based and differential-oracle fuzzers register nothing. We
call this the \emph{oracle gap}.

We present DeFuzz, an agentic system that closes the oracle gap and searches
the matrix directly. First, we systematize the safety invariants that each
defense must satisfy---distilled from compiler source comments, ABI
specifications, and confirmed historical bugs---into a machine-checkable
oracle; every checker returns a four-state verdict and stays sound by grounding
each bug claim in a deterministic binary or runtime signal rather than in model
output. Second, we introduce a cross-mechanism invariant-generation pipeline
that treats a confirmed root cause in one mechanism as a probe, retrieves
structurally analogous code in another mechanism, and instantiates a new
invariant behind two static grounding gates; from 25 confirmed defense bugs and
24 probes over GCC~16.1 it yields 11 new cross-mechanism invariants, none
duplicating an existing seed. Third, we build an explicitly-orchestrated
agentic loop: a deterministic pipeline drives build, coverage, and oracle in a
fixed order and invokes agents only at fixed positions, where they communicate
through a versioned blackboard rather than free-form dialogue. Unlike
free-running agent systems, this design yields reproducible traces, per-edge
ablation, and bug claims that trace back to deterministic evidence; a
checker-bound ISA routing scheme collapses the mechanism $\times$ ISA Cartesian
product into a single semantic decision. On GCC and LLVM, DeFuzz discovered
\tbd{N} silent-failure defects, of which \tbd{M} were confirmed upstream.
\end{abstract}

\section{Introduction}\label{sec:intro}

The number of publicly disclosed vulnerabilities has grown for a decade; the
CVE database now registers more than twenty thousand new entries per year and
the rate is still rising. Major vendors have conceded that manual auditing and
secure-coding disciplines alone are insufficient to eliminate memory-safety bugs in
upper-layer code. As a consequence, the behavior of a program following a bug trigger 
has become increasingly critical: stack canaries, \texttt{\_FORTIFY\_SOURCE},
Indirect Branch Tracking (IBT), Control-Flow Integrity (CFI), Pointer
Authentication (PAC), and shadow stacks are deployed jointly by the compiler
and the runtime as a backstop intended to hold even when an exploit fires.
\incfig{defense-role}{Compiler-inserted defenses as the runtime last
line: they must hold once an upper-layer memory-safety bug is triggered.}

Compiler-emitted defenses share a common working model: at compile time the
compiler inserts check instructions, layout constraints, or metadata so that
any run-time deviation from the intended behavior is caught and the program is
terminated. The stack canary is the canonical example---a random sentinel is
placed between the buffer and the saved return address, and is compared against
its original value before the function returns.

However, defense mechanisms themselves consist of code, which is susceptible to bugs. 
Furthermore, the effectiveness of a defense depends not only on the correctness 
of its implementation but also on the underlying ISA. As the number of ISAs keeps growing, the
cross-architecture attack surface widens, and the rigor of defense implementations 
within compilers becomes correspondingly harder to guarantee.

\subsection{Silent failure}
CVE-2023-4039 is a representative case. The \texttt{-fstack-protector} mechanism of GCC was
defeated on AArch64 for a long time: when a function uses a variable-length
array (VLA) or \texttt{alloca}, the compiler places the canary \emph{above} the
VLA while the saved return address sits \emph{below} it, so an overflow rewrites
the return address before the canary is ever checked. The bug had existed since
early GCC versions and was only analyzed and fixed by outside researchers in
2023. Figure~\ref{fig:stack-layout} shows the buggy frame layout.
\incfigh{stack-layout}{CVE-2023-4039 buggy stack frame: canary above the
VLA, return address below it, so an overflow bypasses the check.}{0.86}

This class of failure has a \emph{double-silent} character: it raises no error
at compile time (toolchain and functional tests see nothing) and produces
correct functional results at run time (program behavior sees nothing), yet the
security contract is already broken. We call this a \emph{silent failure}: the
compiled artifact appears to carry a complete line of defense, but an attacker
can bypass it entirely.

\subsection{Two coupled challenges}
Systematically exposing silent failures presents two intertwined challenges.

\emph{Large search space.} As shown in Figure~\ref{fig:defense-matrix},
the compiler must uphold a security contract across a two-dimensional space of
many defense mechanisms (canary, \texttt{\_FORTIFY\_SOURCE}, IBT, CFI, \dots)
and many ISAs (x86-64, AArch64, RISC-V, LoongArch, \dots). Each cell maps to an
independent middle-end pass or backend template, and any detail can decide
whether the mechanism takes effect. Homologous omissions are not rare: PR-96191,
a sibling of CVE-2023-4039 fixed in 2020, covered only some architectures and
left several fallback backends untouched, so the same failure persisted
silently on other targets for years. Our preliminary experiments already
confirmed several such homologous failures with the GCC upstream.
\incfig{defense-matrix}{The defense mechanism $\times$ ISA matrix. Each
cell is an independent backend/pass; silent failures hide in individual cells.}

\emph{Difficulty in failure adjudication: the oracle gap.} Even when a test reaches
the relevant code path, ``is the defense contract still satisfied'' is not
directly observable: the program does not crash, its output agrees with other
compilers, and everything looks normal while the contract has already been
violated. We call the static or dynamic properties a defense must satisfy at
run time its \emph{safety invariants} (e.g., the canary must lie between any
overflow-reachable stack object and the return address); breaking them is a
silent failure. Existing automated methods are oblivious to this phenomenon: crash- or
differential-output tools such as Csmith and YARPGen trigger no anomaly, and
coverage-guided fuzzing lacks the oracle to judge the security status of each
cell.

\subsection{Our approach}
DeFuzz combines three stages. First, it generates safety invariants through two
complementary paths: a segmented chain-of-thought (CoT) review extracts
candidates from defense-related corpora, while a RAG pass uses confirmed
historical bugs as probes to find implementations with homologous root causes
across mechanisms and compiler backends and generalize them into additional
invariants. Second, accepted invariants are translated into static or dynamic
checkers. Third, the
orchestrator assigns one invariant--checker pair at a time, expands the pair to
the ISAs declared by the checker, and runs an oracle-grounded agent loop. The
model proposes seeds and feedback; deterministic binary or runtime evidence
decides whether a bug exists.

\subsection{Contributions}
\begin{itemize}
  \item \textbf{An oracle for silent failure.} A sound,
  machine-checkable oracle that fills the silent-failure oracle gap and grounds
  confirmed bug reports in reproducible evidence
  (\S\ref{sec:oracle}).
  \item \textbf{Systematized security invariants.} A bottom-up
  taxonomy over 468 catalogued invariants across more than two dozen defenses,
  each in a machine-decidable static/dynamic form. Segmented CoT review provides
  broad corpus coverage, while historical-bug RAG yields 11 new invariants
  (\S\ref{sec:invariants},~\S\ref{sec:specgen}).
  \item \textbf{An explicitly-orchestrated agent system.} A deterministic
  pipeline with agents at fixed positions linked through a versioned blackboard,
  yielding reproducible, ablatable, and auditable runs---in contrast to
  free-running agent systems (\S\ref{sec:loop}).
  \item \textbf{An agentic loop.} Oracle-grounded, coverage-guided
  seed generation scheduled by invariant--checker pairs, with ISA applicability
  derived from checker metadata (\S\ref{sec:loop}).
\end{itemize}

\section{Background and Motivation}\label{sec:bg}

\subsection{Compiler defenses and the mechanism $\times$ ISA matrix}
A compiler-emitted defense inserts, at compile time, either extra check
instructions (canary comparison, CFI target checks), layout constraints (guard
placement, shadow-stack slots), or metadata (BTI landing pads, PAC signing) so
that a run-time deviation is detected and the program is halted. Whether the
inserted defense is actually effective is decided on a per-backend basis: stack
layout, register allocation, and instruction encoding all differ across ISAs,
and each is realized by a distinct pass or target template. The matrix in
Figure~\ref{fig:defense-matrix} therefore has genuinely independent cells, and
a defense that is correct in one cell can be silently broken in the next.

\subsection{Silent failure: a real case}
Returning to CVE-2023-4039: the effectiveness of the stack protector depends on
the invariant ``the canary lies between every overflow-reachable stack object
and the saved return address.'' On AArch64 with a VLA, that invariant is
violated by construction, yet nothing in the build or in a functional test
reveals it. PR-96191 demonstrates a recurrence of this pattern: a fix that lands on some
targets but not on fallback backends leaves the invariant violated elsewhere.
Both cases exemplify the same phenomenon---a complete-looking defense that an attacker
bypasses.

\subsection{Why existing methods miss it}
\emph{Static analysis} is constrained by patterns: silent failures exhibit heterogeneous
forms, and most confirmed cases involve idiosyncratic violations, rendering pattern
matching difficult to generalize. \emph{Crash and differential fuzzing}
monitor surface-level signals---whether the program crashes or disagrees with another
compiler---which by construction cannot catch a silent logic failure; this is
the oracle gap. \emph{Free-running agent fuzzers} can find bugs but sacrifice
reproducibility: with an agent deciding the control flow, when to compile, and
when to declare a bug, trajectories cannot be replayed, ablations cannot be
run, and a bug claim cannot be traced back to a deterministic cause.

\subsection{Scope and threat model}
We target the compiler's own defense implementation on GCC and LLVM. A finding
is a \emph{silent failure}: a compiled artifact whose defense contract is
violated while it compiles cleanly and runs correctly. Following our project
constraint, we describe only the violated invariant (the falsification surface)
and do not construct working exploits. The source-comment and RAG corpus is
pinned to GCC~16.1 to keep evidence and reproduction anchored to one tree.

\section{System Overview}\label{sec:overview}

DeFuzz separates offline knowledge construction from online bug search. The
offline stage reads compiler source, specifications, and historical bug records.
A segmented CoT review provides broad candidate extraction, while historical-bug
RAG targets implementations that may share a known root cause. Candidates that
survive evidence grounding and deduplication enter the invariant catalog and are
implemented as static or dynamic checkers. Confirmed findings can be added to
the historical-bug pool as new retrieval probes.

The online stage tests the compiler rather than the generated program. At run
initialization, the orchestrator enumerates the checker catalog into a fixed
queue. Each target denotes one invariant--checker pair; checker metadata supplies
the applicable ISAs. The target remains fixed until its round or wall-clock budget is exhausted.
Within each round, the Generator produces a C seed and a hard-wired
pipeline executes routing $\rightarrow$ build $\rightarrow$ coverage
$\rightarrow$ oracle. Pass, NotApplicable, and Error enter the Feedback branch
before the next round; Error remains visible as an infrastructure diagnostic
and consumes the current round without becoming a bug. A Fail invokes the
Minimizer and is rechecked against the same checker. All intermediate states and
results are versioned for replay and audit.
\wincfig{architecture}{DeFuzz technical route: corpus-grounded invariant
generation feeds programmatic checkers, which define the targets of an
explicitly orchestrated agent loop.}

Figure~\ref{fig:main-loop} shows one target-specific iteration end to end; the
following sections detail each stage.
\incfigh{main-loop}{Agent loop for one invariant--checker target and its
applicable ISAs: generate $\rightarrow$ build $\rightarrow$ coverage
$\rightarrow$ check $\rightarrow$ feedback or minimization.}{0.92}

\section{Security Invariants for Compiler Defenses}\label{sec:invariants}

\subsection{What is a security invariant}
A \emph{security invariant} is a property a defense must satisfy for it to be
effective; violating it means the mechanism has been silently weakened in the
compiled artifact or at run time, whether or not the code itself contains an
overt bug. Each invariant is recorded with a machine-verifiable
\texttt{statement} and, separately, an \texttt{observation} field that names
only the externally observable phenomenon of a violation---not how to detect it.
The separation of the phenomenon from its detection recipe is deliberate: ``the
\texttt{\_\_stack\_chk\_fail} symbol is missing from the binary'' is a
phenomenon, whereas ``look it up in the ELF \texttt{.dynsym}'' is a detection
mechanism that belongs to the oracle implementation.

\subsection{Sources and segmented coarse extraction}
The invariant archive draws on three source classes: (i) compiler source
comments and assertions, (ii) compiler and ABI documentation, and (iii)
confirmed historical bugs and their patches. Every record carries a uniform
schema---\texttt{ID}, \texttt{statement}, \texttt{compiler}, \texttt{version},
\texttt{target}, \texttt{source\_kind}, evidence, \texttt{version\_sensitivity},
and \texttt{observation}---so that a downstream consumer (oracle, static
analysis, CI) can decide how to use each invariant independently.

The corpus is too large for a single model context. Related evidence also lies
in distant source files, specifications, and patches. We therefore partition
the material by defense mechanism, compiler phase, and function boundary. A segmented CoT
prompt~\cite{wei2022chain} asks the model to identify the protected asset, activation conditions,
required security property, and externally observable violation in each segment.
This pass produces candidates rather than accepted invariants. Every candidate
must retain its source evidence and pass the grounding checks in
\S\ref{sec:specgen}.

\subsection{A bottom-up taxonomy}
Rather than imposing a top-down category scheme, we cluster the 468 catalogued
invariants data-first to avoid preset bias, and read the root-cause families off
the clusters (Figure~\ref{fig:invariant-taxonomy}). Three families recur across
mechanisms and ISAs: (a) \emph{exit-time sensitive-register/state residue}, in
which a fast or exceptional exit bypasses the register/state clearing that a
defense relies on; (b) \emph{stack-clash frame-size truncation}, in which a
forced integer narrowing flips a very large frame size negative and the probing
loop is skipped; and (c) \emph{RISC-V large-address materialization}, in which
an unintended sign extension while splitting or synthesizing a large
address/offset shifts the computed address.
\incfig{invariant-taxonomy}{Bottom-up root-cause families over 468
invariants; three families recur across mechanism and ISA.}

\subsection{Machine-checkable form: static vs dynamic}
Each invariant is classified by whether its violation can be decided from the
un-executed binary (symbols, sections, disassembly, strings)---\emph{static}---or
only by running it and observing exit status, output flags, or register/memory
snapshots---\emph{dynamic}. The choice is guided by four dimensions:
observability, decisiveness, cost, and static/dynamic attribution. When an
invariant has both a statically observable part and a more precise dynamic part,
we split it into two checkers rather than forcing one classification.

\section{Corpus-Grounded Invariant Generation}\label{sec:specgen}

\subsection{Why direct review needs retrieval}
Segmented review gives the LLM broad corpus coverage, but it does not reliably
cluster fragments that encode the same root cause. The relevant evidence may be
separated by mechanism names, compiler phases, target backends, or years of
patch history. RAG~\cite{lewis2020retrieval} therefore provides a complementary
high-yield path: a confirmed historical bug becomes an explicit probe for
related implementations rather than an association the model must discover
while reading unrelated segments. RAG raises the chance of recovering bugs that
share a known root cause; it does not guarantee completeness. Segmented review
covers the remaining corpus for invariants with no historical probe.

\subsection{Historical-bug RAG and root-cause generalization}
The RAG stage first distills a confirmed bug into a mechanism-agnostic failure
signature. Retrieval then searches the indexed middle-end and target-backend
corpus for structurally or semantically related implementations. We avoid
entity-only retrieval based on API or symbol names because such queries tend to
remain within the original mechanism. Instead, the query represents the
abstract operation, violated assumption, and protected asset. This formulation
allows one bug to expose analogous sites across mechanism and backend
boundaries.

\subsection{Pipeline: distill $\rightarrow$ analogy $\rightarrow$ specialize $\rightarrow$ entailment}
A confirmed root cause is first \emph{distilled} into a mechanism-agnostic
description. The \emph{analogy} step decides whether a retrieved code block
instantiates the same failure shape. The \emph{specialize} step instantiates the
invariant on the concrete target source, and \emph{entailment} checks that the
statement follows from the cited evidence. The analogy gate rejects lexical
matches that do not share the root-cause structure before specialization.
Figure~\ref{fig:specgen-pipeline} shows how segmented review and the RAG path
feed a shared grounding and novelty filter.
\incfigh{specgen-pipeline}{Corpus-grounded invariant generation: segmented CoT
review and historical-bug RAG feed a shared grounding and novelty filter;
confirmed new findings return as retrieval probes.}{0.82}

\subsection{Retrieval backends and current yield}
Within this single RAG stage, we evaluate sparse BM25 and dense
\texttt{doubao-embedding-vision} retrieval over the same 24 probe seeds and the
same 4{,}496 GCC-16.1 source chunks. BM25 excels on distinctive identifiers
(\texttt{ix86\_zero\_call\_used\_regs}, \texttt{SUBREG\_PROMOTED},
\texttt{CONST\_HIGH\_PART}); the dense retriever pulls in semantically
isomorphic blocks whose lexical anchors are weak (the CMSE clear-set constructor,
the RISC-V saturating truncation). Deduplicating by \texttt{(seed\_id,
chunk\_id)} and statement, the two paths yield 5 in the intersection, 4
BM25-only, and 2 embedding-only, for 11 in the union
(Figure~\ref{fig:bm25-vs-embedding}); all 11 are marked \texttt{is\_novel} after
a novelty check (threshold 85.0). Table~\ref{tab:retrieval} summarizes.
\incfig{bm25-vs-embedding}{Union/dedup of BM25 and embedding retrieval:
5 intersection, 4 BM25-only, 2 embedding-only, 11 union.}

\begin{table}[t]
\centering
\small
\caption{Cross-mechanism invariants by retrieval backend (24 probes, 4{,}496
GCC-16.1 chunks).}
\label{tab:retrieval}
\begin{tabular}{@{}p{0.37\columnwidth}cp{0.43\columnwidth}@{}}
\toprule
Category & Count & Note \\
\midrule
Intersection (BM25 $\cap$ embed) & 5  & same seed \& chunk, isomorphic statement \\
BM25-only                        & 4  & embed missed the block for that seed \\
Embedding-only                   & 2  & BM25 missed; dense retrieval added \\
\midrule
\textbf{Union (deduplicated)}    & \textbf{11} & $9 + 7 - 5$; all \texttt{is\_novel} \\
\bottomrule
\end{tabular}
\end{table}

\subsection{Grounding gates and reproducibility}
Candidates from both segmented CoT review and historical-bug RAG enter the same
acceptance path. An invariant is accepted only after two static grounding gates
confirm the cited evidence supports its statement and observation, and after the
novelty check (threshold 85.0) confirms it does not duplicate an existing seed
or manual rule.
Because the dense retriever returns slightly different vectors per call, we cache
query vectors keyed by query text so that top-$k$ ordering---and hence the
accepted set---reproduces across runs. After a new finding, including a zero-day,
is reproduced and its root cause is confirmed, we serialize it in the same
schema and add it to the probe pool. The next campaign can then search for
further instances of that root cause.

\section{Oracle: From Invariants to Verdicts}\label{sec:oracle}

\subsection{From an invariant to a checker}
An accepted invariant is not yet an oracle. We first derive a checker contract
that fixes the preconditions, observable signal, applicable ISAs, and evidence
required for a verdict. If the violation is decidable from an unexecuted binary,
the invariant becomes a static checker over symbols, sections, or disassembly.
If adjudication requires exit status, output, or register and memory state, it
becomes a dynamic checker. When both views are useful, we keep two checkers
rather than weakening a deterministic claim into a heuristic.

Each checker returns a four-state verdict
(pass / fail / not-applicable / error). \texttt{NotApplicable} is returned
when the mechanism is off or the seed does not exercise it, and is
ignored by the aggregation layer. A finding is counted as confirmed only when a
deterministic checker fails or an execution differential reproduces, never on
the strength of model output alone
(Figure~\ref{fig:oracle-flow}).
\incfigh{oracle-flow}{An accepted invariant is translated into a static or
dynamic checker; four-state verdicts pass through a deterministic-evidence gate
before any finding is confirmed.}{0.82}

When a programmatic checker cannot decide an invariant--ISA cell and returns
NotApplicable or Error, an optional LLM fallback can inspect extracted
disassembly and symbol evidence. A high-confidence model Fail may promote the
run's provisional aggregate to violated and send the seed to minimization. The
result remains explicitly marked as LLM-derived and is excluded from confirmed
findings until a deterministic checker or execution differential reproduces it.

\subsection{Mechanism aggregator and flag profiles}
A per-mechanism aggregator combines checker verdicts: any Fail raises a
mechanism-level violation carrying its evidence. On/off flag profiles let the
oracle adjudicate differentially (defense enabled vs.\ disabled), and because the
checkers key on binary and runtime signals rather than on the producing
compiler, GCC and LLVM share the same checkers.

\subsection{Checker metadata as a single source of truth}
Each checker declares three metadata fields: which ISAs it binds, whether it is
single-ISA or differential, and whether it is cheap or expensive. This metadata
is the single source of truth read by both the gRPC and the MCP adapter, and it
is what makes the ISA routing in \S\ref{sec:loop} a table lookup rather than an
agent decision.

\section{Agentic Loop}\label{sec:loop}

\subsection{Design principles}
The orchestrator, rather than an agent, owns target selection and control flow.
At run initialization it constructs a fixed queue of invariant--checker pairs.
For each target, checker metadata determines the applicable ISAs, and the target
remains fixed for a configured round or time budget. Within a round, the
pipeline follows the same order: generate $\rightarrow$ routing $\rightarrow$
build $\rightarrow$ coverage $\rightarrow$ oracle $\rightarrow$ route. Agents
are invoked only at fixed positions and communicate through shared state. The
Generator proposes; the oracle stack adjudicates, and confirmed findings require
deterministic evidence.

\subsection{Explicit orchestration and the blackboard}
The agents form a closed loop: feedback guides the next Generator round, the
oracle adjudicates, a non-violation returns to feedback, and a violation goes to
the Minimizer. All communication flows through a versioned blackboard rather
than direct agent-to-agent messages
(Figure~\ref{fig:blackboard}). The blackboard holds the current
invariant--checker target, its applicable ISAs and budget, the seed corpus and
lineage, cumulative coverage, oracle verdicts, feedback guidance, and reduction
results. This design provides three properties: \emph{reproducibility} (pin a blackboard version and any agent's input is
fully determined), \emph{ablatability} (each linkage edge can be switched off
independently), and \emph{auditability} (a bug claim traces back through the
blackboard to a deterministic cause).
\incfig{blackboard}{The blackboard: all inter-agent linkage flows through
versioned shared state, enabling reproduce / ablate / audit.}

\subsection{Three agents}
The \textbf{Generator} (invoked at the start) carries read-only source-search and
invariant-query tools and outputs a seed for the assigned invariant--checker pair;
coverage is not an agent-reportable tool but is measured by the coverage node and
fed back, so the agent cannot fabricate coverage. The \textbf{feedback agent}
(invoked on a non-violation) is a context-isolated subagent that turns the
coverage delta, checker identity, and the Pass/NotApplicable verdict into
next-round guidance
written back to the blackboard. The \textbf{minimizer} (invoked on a violation)
shrinks the PoC with deterministic delta debugging (C-Reduce) as the workhorse
and the LLM only for semantic guidance, so that reduction does not accidentally
delete the structure that triggers the bug. Each candidate reduction is compiled
and rechecked with the same checker before it is retained.

\subsection{Invariant--checker scheduling and ISA binding}
Each target in the deterministic queue denotes one invariant and its checker.
The routing node expands that checker to every ISA in its metadata; the Generator
never chooses either the target invariant or an ISA
(Figure~\ref{fig:checker-routing}). A differential checker always expands to its
full ISA set so that a cross-ISA signal cannot be pruned by model behavior. Cheap
static checkers remain enabled as a side path, but they do not change the
Generator's assigned target. When a checker returns Pass or NotApplicable, the
Feedback agent records the relevant context and the same pair continues until
its budget expires. Error follows the same non-bug path, remains visible as an
infrastructure diagnostic, and consumes the current round. A Fail transfers the
seed and deterministic evidence to the Minimizer. The blackboard records both
branches before the orchestrator advances.
\incfig{checker-routing}{Deterministic target scheduling: one
invariant--checker pair remains fixed for its budget, while checker metadata
expands the applicable ISAs and cheap static checkers remain enabled.}

\section{Implementation}\label{sec:impl}
The Go core (\texttt{cmd/defuzz-core}) exposes \texttt{grpc\_server.go} (build,
coverage, oracle, and checker-metadata services) and \texttt{mcp\_server.go}
(the read-only tools \texttt{search\_source}, \texttt{query\_invariants},
\texttt{coverage\_diff}, \texttt{creduce\_run}, \texttt{compile\_exec}); the
oracle, checkers, and metadata SSOT live under \texttt{internal/oracle}. The
Python orchestrator (\texttt{defuzz\_loop}) holds \texttt{graph.py} (fixed edge
order + conditional routing + enumeration cursor), \texttt{state.py} (the
blackboard), \texttt{audit.py} (per-run directory + checkpointer),
\texttt{routing.py} (the checker~$\rightarrow$~ISA matrix expansion), the three
agents, and the deterministic node wrappers. Each execution generates a dedicated per-run
directory with \texttt{checkpoints.sqlite} (the versioned state chain) and
\texttt{manifest.json} (git SHA, toolchains, checker catalog, LLM config,
ablation flags), which the \texttt{inspect} / \texttt{replay} / \texttt{trace-bug}
subcommands consume for reproduction.
The target queue stores checker IDs, each of which names the corresponding
invariant--checker pair; the cursor advances only after that target's round or
time budget is exhausted.

\section{Evaluation}\label{sec:eval}
\emph{Setup.} Targets are an instrumented GCC~16.1, a cross-GCC, and LLVM;
ISAs are x86-64, AArch64, RISC-V, and LoongArch, with cross-architecture
execution under QEMU user mode. We answer five research questions.

\paragraph{RQ1 --- Real bugs} DeFuzz found \tbd{N} silent-failure defects on
GCC/LLVM; \tbd{M} were confirmed upstream (\tbd{CVE IDs}).

\paragraph{RQ2 --- Invariant generation quality} The historical-bug RAG path
produced 11 union invariants (5 intersection, 4 BM25-only, 2 embedding-only),
all \texttt{is\_novel}; the analogy gate filtered roughly 94\% of BM25 hits as
lexical collisions (Table~\ref{tab:retrieval}).

\paragraph{RQ3 --- Ablation} We toggle coverage feedback, oracle grounding, and
each inter-agent edge, and compare checker routing against a full Cartesian
fan-out (Figure~\ref{fig:ablation}). \tbd{ablation numbers}
\figplaceholder{ablation}{Ablation: contribution of oracle grounding, coverage
feedback, each inter-agent edge, and checker routing to hit rate/cost.}

\paragraph{RQ4 --- Explicit orchestration} Pinning a blackboard version
reproduces a run's trajectory exactly; we compare stability and hit rate against
free orchestration. \tbd{reproducibility numbers}

\paragraph{RQ5 --- Agentic vs.\ fuzz loop} We measure the gain on ``reaching
silent-failure bugs'' over the prior coverage-driven fuzz loop.
\tbd{comparison numbers}

\section{Discussion and Limitations}\label{sec:disc}
DeFuzz describes only violated invariants and does not build exploits; the RAG
corpus is pinned to GCC~16.1; dense retrieval is non-deterministic and requires
query-vector caching to reproduce; and checker routing is a performance
optimization guarded by always-on static checkers, not a correctness mechanism.

\section{Related Work}\label{sec:related}
\emph{Compiler bug finding.} Csmith~\cite{yang2011csmith},
YARPGen~\cite{livinskii2020yarpgen}, and EMI~\cite{le2014emi} find miscompilations
with crash or differential oracles, which by construction miss silent \emph{defense}
failures. \emph{Silent/security bugs in compilers.} Studies such as ``Silent Bugs
Matter''~\cite{silentbugs2023} taxonomize compiler-introduced security bugs but do
not build an agentic detector. \emph{LLM/RAG property generation.}
PropertyGPT~\cite{liu2025propertygpt} generates and verifies smart-contract
properties via retrieval-augmented generation; SpecAuditor~\cite{lin2026specauditor}
detects specification violations by semantic generalization---we deliberately
reject its entity-driven retrieval in favor of abstract-failure-mode transfer.
\emph{LLM-agent fuzzing.} AgentFuzz~\cite{liu2025agentfuzz} applies directed
greybox fuzzing to LLM agents; in contrast, our agents are explicitly
orchestrated rather than free-running.

\section{Conclusion}\label{sec:concl}
An invariant-grounded oracle turns ``agents can find bugs'' from a one-off demo
into a reproducible method for silent defense failures across the mechanism
$\times$ ISA matrix. By pairing a sound oracle with an explicitly-orchestrated
agentic loop and checker-bound ISA routing, DeFuzz makes such bugs both findable
and auditable.

\bibliographystyle{IEEEtran}
\bibliography{refs}

\end{document}
